\section{Idea Density}

\textbf{Idea Density} = Number of propositions/number of words.

A proposition contains information and can be true or false.

Idea density can be used as a measure of intelligence and a low density may indicate some mental impairment. However, to be truly useful, that quality of the ideas must also be taken into account. First and foremost is specificity. A specific idea is of more value than a vague one. We recently suggested that in \textit{Spirituality for the XXIst Century}\footnote{\url{https://gornahoor.net/?p=8630}} an idea that can be manifested can be of special value.

As an example, the post \textit{Salvation and Evolution}\footnote{\url{https://gornahoor.net/?p=8642}}, particularly Tomberg's list, shows high density. It defies summary, which must necessarily be more dense, yet invites expansion. There are specific names and movements that can be discussed. 

\paragraph{Intellectual Pride}
There is often a misunderstanding by religious people concerning esoterism. \textbf{Frithjof Schuon} writes:

\begin{quotex}
There are people who are quite ready to attribute to intelligence a Luciferian aspect and who speak without hesitation of ``\textbf{in\-tel\-lec\-tu\-al pride}" as if this were not a contradiction in terms; hence also the exaltation of ``childlike" or ``simple" faith, which indeed we are the first to respect when it is spontaneous and natural, but not when it is theoretical and affected. 

\end{quotex}
Schuon responds to that charge:

\begin{quotex}
It is not the man who chooses the way. It is the way that chooses the man. In other words. the question of a choice does not arise, since the finite cannot choose the Infinite: rather the question is one of vocation, and those who are ``called" to use the Gospel expression, cannot ignore the call without committing a ``sin against the Holy Ghost" any more than a man can legitimately ignore the obligations of his religion. 

\end{quotex}
Gnosis is one of the gifts of the Holy Spirit, which therefore is not to be spurned. So, on the contrary, intellectual pride can be ascribed to those who claim a knowledge that they are not entitled to, which can only be a simulacrum of true gnosis.

\paragraph{The Scent of a Woman}
Frithjof Schuon points out that, according to an Hadith, \textbf{women, perfumes, and prayer} are worthy of love. Schuon explains why:

\begin{itemize}
\item \textbf{Woman}, synthesizing in her substance virgin nature, the sanctuary and spiritual company, is for man what is most lovable; in her highest aspect, she is the formal projection of merciful and infinite Inwardness in the outward; and in this regard she assumes a quasi-sacramental and liberating function. 
\item \textbf{Perfumes} represent qualities or beauties that are formless, exactly in the same way as music; that is to say that side by side with the formal projection of Inwardness, there exists also a complementary formless projection, symbolized, not by visual or tangible qualities, but by auditory and olfactory ones; perfumes are silent music. 
\item \textbf{Prayer} leads from Outward to Inward, and both consecrates and transmutes the qualitative elements of the outward realm. 
\end{itemize}
We are accustomed to associate Love with the Goodness of God, but Mystical Love relates to the Beauty of God. Schuon even claims that the Beauty of God corresponds to a deeper reality than His Goodness. No doubt, Prayer may take some effort and training in order to learn to love it, but we instinctively have fond memories of scents and aromas from the past – a fresh baked pie, newly mown grass, or the passing whiff of a good perfume.

Those men who have no experience of the attraction to a woman are really missing out – and not a marginal or arbitrary experience, but an essential one. Poetry and literature abound with stories of that attraction, which is beyond rational explanation. Hence, we can suffice with a more homely example.

Some time ago, ``Sophia" was a model who took a job as a project manager at a tech company. Although there were other attractive woman in the group, Stella stood out in an inexplicable way. Perhaps it was her stylish clothes, her smile, or just the way she moved. Well, at 3 PM every afternoon, Stella would go outside for a walk. Soon, large groups of men would start taking their ``breaks" outside at the same hour.

\paragraph{Two Heavens}
\begin{quotex}
There are two heavens: one sensible, and the other, ideal. In the one, both the animal spirits and the rational souls enjoy bliss; in the other, the rational souls alone. The latter paradise is the heaven of knowledge and intuition. \flright{\textsc{Ibn Arabi}}

\end{quotex}
He goes on to explain:

\begin{quotex}
God has depicted paradise in accordance with the different degrees of man's understanding. The Messiah [Christ] defined the delights of paradise as purely spiritual … if the Messiah was so explicit on this point and had recourse to none of the allegories found in our Book, it was simply because his words were spoken to a people … whose mind was thus prepared for his words.

Not so with our Prophet Mohammed whose Divine mission fell among a rude people … Accordingly, most of the descriptions of paradise in his book are based on the body, in order that they might be understood by the people and serve as an incentive to their minds. 

\end{quotex}
Interesting idea, but one can hardly count on ``rational" souls (i.e., pneumatics) being prevalent today. I suspect there are many who get more excited about the prospect of uniting with their pet dogs in the afterlife, rather than experiencing the Beatific Vision.

\paragraph{Die Before you Die}
A common criticism is that the prospect of a future postmortem life devalues human life in the here and now. Yet, to be honest, how many of your human projects have come to fruition? Sure, there are precious few who have excelled in wealth, fame, beauty or intelligence, only to find they are ephemeral. If you have not sought those things, perhaps you have just set your sights too low.

However, if we understand the human condition correctly, we will better understand our project and our destiny. If we have fallen from a higher state, then, like trapped miners, we need to find a way to restore that higher state.

For exoteric religion, that is ``salvation" which comes when we are freed from the body by death. The esoteric task, on the other hand, is liberation while in the body. That is ``dying before you die."

\paragraph{Image and Dignity}
There has been a subtle change in the teaching of the Roman Church since Vatican II, which accounts for much of what churchmen today believe and teach. There are three related concepts:

\begin{enumerate}
\item Man is born in the image and likeness of God 
\item We are all children of God 
\item The human person has ipso facto ``dignity" 
\end{enumerate}
The more Traditional view is that (1) applies to Adam, but the rest of mankind has lost the ``likeness". The ``image" of God means to have an intellectual soul as well as a body. The ``likeness" means to be deified (or, in the West, to receive ``sanctifying grace"). That comes from the second birth.

Hence, only those who have the likeness can become adopted children of God, that is, to be of the same nature, a ``chip off the old block". Finally, dignity is ``the state or quality of being worthy of honor or respect", so the true dignity of man consists in participating in the divine nature. (see \textit{Is the New Catechism Catholic?}\footnote{\url{http://sspx.org/en/new-catechism-catholic}} for the Roman teaching and \textit{The Image And Likeness Of God}\footnote{\url{http://www.stgeorgeserbian.us/darren/darren03.htm}} for the Greek).

\paragraph{Toleration in the Empire}
Neo-pagans like to praise the alleged ``toleration" of the pagans (as though that were truly a virtue). However, life in the Ancient City was hardly pluralistic. When religions clashed, one had to become subservient to the other in hierarchical fashion.

Likewise, in the Roman Empire, colonial religions were tolerated as long as they put allegiance to the Emperor in the highest place. The only exception was the Hebrews, since the Romans venerated great antiquity. Nevertheless, that exception was not passed onto the Christians, since they were not considered a Hebrew cult.

In today's world, Liberalism holds the same place as the Emperor. A man is free to practice any religion he pleases, as long as, in the event of any clash, he is subservient to the Liberal State. Hillary Clinton, for example, has the hubris to assert that religions need to change their teaching on abortion and homosexuality.

Liberalism is incompatible with true religion. Actually, a marker of spiritual progress is the extent to which a person purges himself of liberal assumptions.

\paragraph{Social Reign of Christ}
The incursion of Islam into the West should bring to light how much of Tradition has been lost. On the contrary, unfortunately, the temptation is to double down on a bad debt.

The talking heads, even those who self-identify as ``conservative" condemn in Islam characteristics that used to be considered ``sane and normal" in Europe. It would be one thing to oppose one Tradition to another, but quite another to adopt Liberalism instead.

For example, it used to be normative for Catholicism to be the state religion wherever it was dominant. Hence, the state would fall under religious law in its various aspects as natural law, social teachings, and so on. Modest dress was encouraged (and still is at Latin masses), the man was head of the family, prayer was the center of life.

Other religions would be allowed as long as worship was private. Moreover, the priest would serve Mass ``ad orientem", i.e., facing East. Anyone proposing such a program today would be shunned.

\paragraph{Star of Creation}
This is from a facebook ``discussion":

\begin{quotex}
\textbf{Nick}: Why does your standard bear the Jewish star of Moloch/Satan?

\textbf{Gornahoor}: You have quite the imagination, Nick. That is my family coat of arms that goes back to the 14th century … perhaps you should ask them.

\textbf{Nick}: Sure, stars are a part of medieval heraldry, and the six-pointed ones especially in countries, where Jews were given rights and a path to infiltration. In the case of your standard, the star is what is commonly called the star of David, but actually of those two demons I mentioned. Can you find out what it's supposed to represent in your family's coat of arms? I doubt it's well-meant, if it's put on top of the cross, considering the followers of that star are the biggest anti-Christians that walk the earth. 

\end{quotex}

Obviously, it is not a ``Star of David" which consists of two interlaced equilateral triangles.

A similar image is the symbol of the Anahata, or Heart, Chakra. This chakra is connected to the Sun, and also to Christ, since that is where one points when saying the Sign of the Cross.

Also, in Hinduism, the six-pointed star refers to Shakti (pointing down) and Shiva (pointing up), so their union signifies ``creation" through the union of male and female. Likewise, in Christianity, the six-pointed star (found in some Orthodox churches) is called the \textbf{Star of Creation}. Hence, the true symbolism should be obvious.

I used to catalog the various kinds of idiots I encountered in life until I ran out of fingers and toes. However, I'll give ``Nick" a special award: \textbf{The 21st Idiot Award} (and an honorable mention to the person who ``liked" Nick's comment).



\flrightit{Posted on 2016-08-02 by Cologero }

\begin{center}* * *\end{center}

\begin{footnotesize}\begin{sffamily}



\texttt{Mark Citadel on 2016-08-10 at 06:56 said: }

Excellent stuff. I was fascination by the two heavens concept, as it relates to my own thinking on racial differences present between the Arabs who embraced Islam, and the Occidentals who embraced Christianity. Interestingly, though we were far more capable of understanding a heaven of spiritual reward, we also had this aversion to an element of mitigated dualism, that being the gulf between God and man. It's as if our needs in the here and now were more physical, we needed a God that was in some senses human, but in the future, we were more amenable to the abandonment of the merely physical, so long as the spiritual destiny was close to God.

What you say about Islam's incursion is entirely accurate. Since I was once active in the Counter-Jihad sphere, I still get people who misunderstand my present opposition to Islam in the Occident for a simple xeno-aversion, when actually I have a lot of respect for Islam.

The main thing I took away from this article however, is that I definitely need to read more Schuon.


\end{sffamily}\end{footnotesize}
